First read the number $N$.

To make a query, print ``? x'' $(1 \le x \le N)$ and read the answer:

\begin{itemize}
\item \texttt{YES} if the number $x$ is favourite;
\item \texttt{NO} if the number $x$ is not favourite.
\end{itemize}

Each number may be asked about at most once, the order of queries is arbitrary. To give the final answer, print ``! F'', where $F$ is the number of favourite numbers, and terminate the program.

Your program will get the ``Wrong answer'' verdict in case of:

\begin{itemize}
\item violating the output format;
\item querying a number outside the range $[1; N]$;
\item querying the same number twice;
\item a wrong final answer.
\end{itemize}

\textbf{After each printed query be sure to flush the output buffer}, otherwise your program will get the ``Idleness Limit Exceeded'' verdict:

\begin{itemize}
\item C++: \texttt{std::endl} when printing via \texttt{iostream} (\texttt{cout << ... << endl}), \texttt{fflush(stdout)} after printing via \texttt{stdio} (\texttt{printf});
\item Java: \texttt{System.out.println(...)}, and when printing via another stream (\texttt{PrintWriter}, \texttt{BufferedWriter}) --- its \texttt{flush()} after each line;
\item Python: \texttt{print(..., flush=True)}, and read the subsystem's answers via \texttt{input()}.
\end{itemize}
